%%%  Ukázkový text a dokumentace stylu pro text závěrečné (bakalářské a
%%%  diplomové) práce na KI PřF UP v Olomouci
%%%  Copyright (C) 2012 Martin Rotter, <rotter.martinos@gmail.com>
%%%  Copyright (C) 2014 Jan Outrata, <jan.outrata@upol.cz>


%%  Pro získání PDF souboru dokumentu je třeba tento zdrojový text v
%%  LaTeXu přeložit (dvakrát) programem pdfLaTeX.

%%  V případě použití programu BibLaTeX pro tvorbu seznamu literatury
%%  je poté ještě třeba spustit program Biber s parametrem jméno
%%  souboru zdrojového textu bez přípony a následně opět (dvakrát)
%%  přeložit zdrojový text programem pdfLaTeX.

%%  Postup získání Postscriptového souboru je popsán v dokumentaci.


%%  Třída dokumentu implementující styl pro závěrečnou práci. Vybrané
%%  nepovinné parametry (ostatní v dokumentaci):

%%  'master' pro sazbu diplomové práce, jinak se sází bakalářská práce

%%  'program=kód' pro Váš studijní program/obor (specializaci), kódy
%%  pro diplomovou práci 'infoi' pro Informatiku (Obecná informatika),
%%  'infui' pro Informatiku (Umělá inteligence), 'ainfpst' pro
%%  Aplikovanou informatiku (Počítačové systémy a technologie), 'uinf'
%%  pro Učitelství informatiky pro střední školy, 'binf' pro
%%  Bioinformatiku, 'inf' pro Informatiku (bez specializací) a 'ainf'
%%  pro Aplikovanou informatiku (bez specializací), jinak je výchozí
%%  ainfvs pro Aplikovanou informatiku (Vývoj software), a pro
%%  bakalářskou práci 'infoi' pro Informatiku (Obecná informatika),
%%  'itp' pro Informační technologie v prezenční formě, 'itk' pro
%%  Informační technologie v kombinované formě, 'infv' pro Informatiku
%%  pro vzdělávání, 'binf' pro Bioinfomatiku, 'inf' pro Informatiku
%%  (bez specializací), 'ainfp' pro Aplikovanou informatiku (bez
%%  specializací) v prezenční formě, 'ainfk' pro Aplikovanou
%%  informatiku (bez specializací) v kombinované formě, jinak je
%%  výchozí infpvs pro Informatiku (Programování a vývoj software)

%%  'printversion' pro sazbu verze pro tisk (nebarevné logo a odkazy,
%%  odkazy s uvedením adresy za odkazem, ne odkazy do rejstříku),
%%  jinak verze pro prohlížeč

%%  'biblatex' pro zapnutí podpory pro sazbu bibliografie pomocí
%%  BibLaTeXu, jinak je výchozí sazba v prostředí thebibliography

%%  'language=jazyk' pro jazyk práce, jazyky english pro anglický,
%%  slovak pro slovenský, jinak je výchozí czech pro český

%%  'font=sans' pro bezpatkový font (Iwona Light), jinak je výchozí
%%  serif pro patkový (Latin Modern)

%%  'figures, tables, theorems a sourcecodes' pro sazbu seznamu
%%  obrázků, tabulek, vět a zdrojových kódů, jinak při =false se
%%  nesází (u theorems a sourcecodes výchozí)

\documentclass[
%  master,
%  program=ainfvs,
%  printversion,
  biblatex,
  seconds=true,
  language=slovak,
%  font=sans,
  figures=false,
%  tables=false,
%  theorems,
%  sourcecodes,
  glossaries,
  index
]{kidiplom}

%% Informace pro úvodní strany. V jazyku práce (pokud není v komentáři
%% uvedeno česky) a anglicky. Uveďte všechny, u kterých není v
%% komentáři uvedeno, že jsou volitelné. Při neuvedení se použijí
%% výchozí texty. Text pro jiný než nastavený jazyk práce (nepovinným
%% parametrem language makra \documentclass, výchozí český) se zadává
%% použitím makra s uvedením jazyka jako nepovinného parametru.

%% Název práce, česky a anglicky. Měl by se vysázet na jeden řádek.
\title{Paralelní Programování v Go}
\title[english]{Paralel programing in Go}

%% Volitelný podnázev práce, česky a anglicky. Měl by se vysázet na
%% jeden řádek. Výchozí je prázdný.


%% Jméno autora práce. Makro nemá nepovinný parametr pro uvedení
%% jazyka.
\author{Richard Bebjak}

%% Jméno vedoucího práce (včetně titulů). Makro nemá nepovinný
%% parametr pro uvedení jazyka.
\supervisor{Mgr. Jan Laštovička, Ph.D.}

%% Volitelný rok odevzdání práce. Výchozí je aktuální (kalendářní)
%% rok. Makro nemá nepovinný parametr pro uvedení jazyka.
%\yearofsubmit{\the\year}

%% Anotace práce, včetně anglické (obvykle překlad z jazyka
%% práce). Jeden odstavec!
\annotation{}

\annotation[english]{}

%% Klíčová slova práce, včetně anglických. Oddělená (obvykle) středníkem.
\keywords{}
\keywords[english]{}

%% Volitelná specifikace příloh textu práce, i anglicky. Výchozí je
%% 'elektronická data v systému katedry informatiky / electronic data
%% in system of department of computer science'.
%\supplements{nejlepší software všech dob}
%\supplements[english]{the best software of all times}

%% Volitelné poděkování. Stručné! Výchozí je prázdné. Makro nemá
%% nepovinný parametr pro uvedení jazyka.
\thanks{Děkuji, děkuji, děkuji.}

%% Cesta k souboru s bibliografií pro její sazbu pomocí BibLaTeXu
%% (zvolenou nepovinným parametrem biblatex makra
%% \documentclass). Použijte pouze při této sazbě, ne při (výchozí)
%% sazbě v prostředí thebibliography.
\bibliography{}

%% Další dodatečné styly (balíky) potřebné pro sazbu vlastního textu
%% práce.
\usepackage{lipsum}
\usepackage{longtable}
\usepackage{listings}
\usepackage{biblatex}
\addbibresource{bibliografie.bib}

\begin{document}
%% Sazba úvodních stran -- titulní, s bibliografickými údaji, s
%% anotací a klíčovými slovy, s poděkováním a prohlášením, s obsahem a
%% se seznamy obrázků, tabulek, vět a zdrojových kódů (pokud jejich
%% sazba není vypnutá).
\maketitle

%% Vlastní text závěrečné práce. Pro povinné závěry, před přílohami,
%% použijte prostředí kiconclusions. Povinná je i příloha s obsahem
%% elektronických dat.

%% -------------------------------------------------------------------

\newcommand{\BibLaTeX}{\textsc{Bib}\LaTeX}

% \noindent\textcolor{red}{\LARGE Upozornění: Následující text
%   dokumentace stylu, vyjma přílohy~\ref{sec:ObsahData}, je rozpracovaná
%   a (značně) neúplná verze!!!}


\section{Komunikačné kanály}
Komunikačné kanály predstavujú jeden zo základných mechanizmov synchronizácie a komunikácie medzi goroutines. Sú zabudované referečným typom,
ktorý umožňuje bezpečný prenos dát. Ich návrch vychádza z Communicating Sequential Processes, 
ktorý zdôrazňuje komunikáciu medzi procesmi miesto zdieľania pamäte.\cite{cox2017}

Komunikačné kanály sa deklarujú pomocou kľúčového slova \texttt{chan}
a vždy sa deklarujú aj s určitým typom. To znamená, že kanál môže prenášať hodnoty
iba tohoto typu. 
Komunikačné kanály patria k referečným typom, preto ich deklarácia nevytvára potrebnú pamäťovú štruktúru.
Pamäť sa alokuje pomocou zabudovanej funkcie \texttt{make}. Takto vytvorený komunikačný kanál umožňuje zapisovanie a čítanie hodnôt.
Zápis do komunikačného kanála umožnuje operátor \texttt{<-}, ktorý sa umiestňuje na pravú stranu názvu komunikačného kanála,
pre čítanie sa operátor umiestňuje na ľavú stranu názvu komunikačného kanála.
\begin{lstlisting}[language=Go]
    var channel = make(chan int, 1)

    go func(){
        channel <- 42
    }()

    result := <- channel
    close(channel)
\end{lstlisting}

Dôležitou vlastnosťou komunikačného kanála je jeho blokujúce správanie. Ak goroutine zapisuje do komnikačného kanála, ktorý je plný, je jej proces
blokovaný kým iná goroutine z komunikačného kanála hodnotu neprečíta. Podobne čítanie z prázdneho komunikačného kanála je blokovaný kým nie je do
komunikačného kanála zapísaná nová hodnota.

Komunikačný kanál je možné zavrieť pomocou funkcie \texttt{close}. Uzavretie sygnalizuje príjmateľovi, že do komunikačného kanála nebude viac zapisované.
Pri čítaní z kanála je možné získať druhú návratovú hodnotu typu bool, ktorá indikuje jeho otvorenie.
\subsection{Bufferované a nebufferované komunikačné kanály}
Každý komunikačný kanál môže byť bufferovaný alebo nebufferovaný. Nebufferovaný komunikačný kanál nemá žiadnu kapacitu a každá operácia zápisu musí
byť okamžite spravovaná operáciou čítania.

Bufferovaný komunikačný kanál sa definuje pomocou voliteľného druhého argumentu funkcie make.
Bufferovaný komunikačný kanál obsahuje bugger s definovanou kapacitou, ktorá umožní uložiť určitý počet hodnôt do buffera bez okamžitého čítania.
Bufferovaný komunikačný kanál zablokuje zápis až v momente keď je buffer plný.
\subsection{Jednosmerný komunikačný kanál}
Jazyk Go umožňuje definovanie jednosmerných kanálov, ktoré podporujú iba zápis alebo čítanie. Tieto komunikačné kanály sa používanú ako parametre 
funkcií, aby určili akým spôsobom sa komunikačný kanál v danej časti programu používa. Tento mechanizmus zvyšuje bezpečnosť programu, pretože
kompilátor dokáže zachytiť nesprávne použitie.

\section{Synchronizačné primitíva}
Jazyk Go poskytuje sync package, ktorý obsahuje synchronizačné primitíva. Tieto synchronizačné primitíva sú používane pri nízko úrovňovom
prístupe do pamäti. Rozdiel medzi ostatnými jazykmi a jazykom Go je, že Go poskytuje novú sadu synchronizačných primitív pre súbežnosť nad 
sadov primitív pre synchronizáciu prístupu k pamäti.

\newpage
\subsection{WaitGroup}
\subsection{Select}
\subsection{Go pamäťový model}
\section{Chat server}
\subsection{Funkčné požiadavky}
\subsection{Súbežná archytektúra}
\section{Paralelné algoritmy}
\section{Paralené programovanie}
\subsection{Paralelné a súbežné programovanie}
\subsection{Problémy súbežného programovania}
\section{Jazyk Go}
\subsection{Goroutines}



%%% Argument `joinlists' způsobí zřetězení seznamů obrázků, tabulek,
%%% vět a zdrojových kódů. Není-li použít, všechny seznamy jsou
%%% uvedeny na samostatných stránkách.


%%  'encoding=kódování' pro kódování tohoto a vložených zdrojových
%%  textů v kódování jiném než výchozím utf8




%% v případě tvorby rejstříku přeložit vygenerovaný soubor .idx
%% programem Makeindex a v případě tvorby seznamu zkratek spustit
%% program Makeglossaries s parametrem jméno souboru zdrojového textu
%% bez přípony a následně opět (dvakrát) přeložit zdrojový text
%% programem pdfLaTeX.


%% Závěry práce. V jazyce práce a anglicky. Text pro jiný než
%% nastavený jazyk práce (nepovinným parametrem language makra
%% \documentclass, výchozí český) se zadává použitím makra s uvedením
%% jazyka jako nepovinného parametru.

%% Přílohy obsahu textu práce, za makrem \appendix.
\appendix

\section{První příloha}
Text první přílohy

\section{Druhá příloha}
Text druhé přílohy

%% Obsah elektronických dat. Poslední příloha. Upravte podle vlastní
%% práce!
\section{Obsah elektronických dat} \label{sec:ObsahData}



%% -------------------------------------------------------------------

%% Sazba volitelného seznamu zkratek, za přílohami.
\printglossary

%% Sazba povinné bibliografie, za přílohami (případně i za seznamem
%% zkratek). Při použití BibLaTeXu použijte makro
%% \printbibliography. jinak prostředí thebibliography. Ne obojí!

%% Sazba i v textu necitovaných zdrojů, při použití
%% BibLaTeXu. Volitelné.
\nocite{*}
%% Vlastní sazba bibliografie při použití BibLaTeXu.
\printbibliography

%% Bibliografie, včetně sazby, při NEpoužití BibLaTeXu.
% \begin{thebibliography}{9}
%\bibitem{kniha2} \uppercase{Hawke}, Paul. NanoHttpd: Light-weight HTTP server designed for embedding in other applications. GitHub [online]. 2014-05-12. [cit. 2014-12-06]. Dostupné z: \url{https://github.com/NanoHttpd/nanohttpd}
%
%\bibitem{jeske13} \uppercase{Jeske}, David; \uppercase{Novák}, Josef. Simple HTTP Server in \csharp: Threaded synchronous HTTP Server abstract class, to respond to HTTP requests. CodeProject: For those who code [online]. 2014-05-24. [cit. 2014-12-06]. Dostupné z: \url{http://www.codeproject.com/Articles/137979/Simple-HTTP-Server-in-C}
%
%\bibitem{uzis2012} \uppercase{ÚSTAV ZDRAVOTNICKÝCH INFORMACÍ A STATISTIKY ČR}. Lékaři, zubní lékaři a farmaceuti 2012 [online]. Praha 2, Palackého náměstí 4: Ústav zdravotnických informací a statistiky ČR, 2012 [cit. 2014-12-06]. ISBN 978-80-7472-089-5. Dostupné z: \url{http://www.uzis.cz/publikace/lekari-zubni-lekari-farmaceuti-2012}
% \end{thebibliography}

%% Sazba volitelného rejstříku, za bibliografií.
\printindex

\end{document}

%%% Local Variables:
%%% mode: latex
%%% TeX-master: t
%%% End:
